\caption{\textbf{The auto-nudger is an inefficient Maxwell demon: its bits buy glances, not committed work (H35, with H60).}
(a)~Goal period \#51 (short-pause scaffold), trap-age state space: work bought beyond a random nudger with the same budget against the state information the policy uses. The logged nudger uses 1.04 bits per nudge and buys $+0.50$ [$-0.09$, 1.16] extra active minutes, $\eta_\mathrm{SU}=0.38$ [$-0.08$, 0.68] of what the frontier $V^*(R)$ gives at the same bits; nudging once at the early re-pauses ($k=2$--3) reaches the frontier's end ($\times2.3$ [1.7, 5.2] work per nudge), nudging only deep traps is worse than random; the shaded region lies above the Donsker--Varadhan ceiling.
(b)~First-nudge response by trap age (95\% CI; dashed, shrunk) and where the nudges go: 38\% fall on chains $\ge10$ pauses deep, where a nudge buys 0.5 minutes against 2.2--2.6 at $k=1$--9.
(c)~Efficiency by outcome (round 1b, gray band: random nudging): glances $\eta_\mathrm{SU}=0.88$ [0.26, 0.97]; committed work 0.19 [$-1.39$, 0.70] after placebo differencing (post hoc), with $+0.22$ [$-0.29$, 0.82] commits per hour per nudge.
(d)~H60, idle gates in \#51: extra active calls in 30\,min per nudge under each policy at the logged budget (cross-fitted, 95\% CI). Random gate choice beats the logged nudger ($\times1.46$ [1.22, 1.67]) and once-early beats it ($\times1.55$ [1.25, 1.81]); a fitted index does not beat once-early ($\times1.08$ [$-0.30$, 3.44]; once per trap $\times1.49$ [0.46, 3.42]).}
\label{fig:H35}
